\documentclass[10pt]{article}
\usepackage[a4paper,margin=25mm]{geometry}
\usepackage{amsmath,amssymb,amsthm}
\usepackage[hidelinks]{hyperref}
\newtheorem{theorem}{Theorem}
\hypersetup{hidelinks}
\begin{document}
\title{A shellable bipartite graph with an induced six-cycle}
\author{Chengzhe Gao}\date{}\maketitle
Conjecture 541 asserts an equivalence between shellability of the
independence complex of a bipartite graph and chordality. We give a
counterexample under both ordinary chordality and chordal
bipartiteness.

\paragraph{Construction.}
Let $v_0,\ldots,v_5$ form a six-cycle (indices modulo $6$), and attach
one new leaf $w_i$ to each $v_i$. These are all twelve vertices and
all twelve edges. Give $v_i$ the color $i\bmod2$ and $w_i$ the
opposite color. This is a proper bipartition. The induced subgraph
on the $v_i$ is $C_6$, so the graph is neither chordal nor chordal
bipartite.

\begin{theorem}
The independence complex of this graph is pure and shellable.
\end{theorem}
\begin{proof}
A maximal independent set contains exactly one of $v_i,w_i$ for
every $i$. It cannot contain both because they are adjacent. If it
contained neither, the leaf $w_i$ could be added, contradicting
maximality. Thus its facets are exactly
\[
 F_I=\{v_i:i\in I\}\;\cup\;\{w_i:i\notin I\},
\]
where $I$ is an independent subset of $C_6$. Conversely, each such
set is independent and maximal, because the omitted member of each
pair is adjacent to the chosen one. Every facet therefore has six
vertices. There are $1+6+9+2=18$ of them.

Order these facets by increasing $|I|$, breaking ties arbitrarily.
Consider an earlier facet $F_J$ and a later facet $F_I$. They are
distinct and $|J|\leq|I|$, so there is an $i\in I\setminus J$.
The independent set $I'=I\setminus\{i\}$ has smaller size and
therefore its facet occurs earlier. Moreover,
\[
 F_I\setminus F_{I'}=\{v_i\}\subseteq F_I\setminus F_J.
\]
This is the standard facet-difference shelling criterion: for every
earlier $F_J$ there is an earlier $F_{I'}$ whose difference from the
current facet is one vertex belonging to $F_I\setminus F_J$.
See \cite{matching}, Definition 2.7. Hence the displayed order is
a shelling. Since the complex is pure, this argument works for the
usual pure definition as well as the general nonpure definition.
\end{proof}

The graph consequently satisfies the bipartite hypothesis and the
shellability side of the proposed equivalence, while failing either
meaning of chordality relevant to the statement.

\paragraph{Formal verification.}
Lean uses vertices $0,\ldots,11$, with $i$ representing $v_i$ and
$i+6$ representing $w_i$. The actual adjacency function is checked
to be irreflexive and symmetric. A concrete two-coloring and an
injective, induced embedding of $C_6$ are checked, giving the two
failures of chordality from their definitions.

Subsets are encoded as twelve-bit masks. The file proves that every
Boolean subset and, classically, every proposition-valued subset
has an encoding. It proves that the computable facet criterion
(independence and inability to add an outside vertex) is equivalent
to genuine inclusion-maximal independence against every subset.
All $4096$ masks are checked, in 64 blocks, to show that the list
of 18 facets is complete and has no repetitions. Every facet is
checked to have six vertices. Finally the exact facet-difference
criterion used in the proof is verified for every ordered pair of
facets. Thus no facet, subset, or required pair is omitted.

The proof uses only \texttt{Std} and Lean 4.19.0. Finite checks are
kernel-reduced \texttt{decide} proofs, with no \texttt{native\_decide},
unproved mathematical axioms, or admitted steps. The final theorem
\texttt{conjecture541\_\allowbreak{}false} negates the asserted characterization;
the stronger counterexample theorem also excludes chordal
bipartiteness.

\begin{thebibliography}{1}
\bibitem{matching} S. Morey, E. Reyes and R. H. Villarreal,
\emph{Cohen--Macaulay, shellable and unmixed clutters with a perfect
matching of K\"onig type}, J. Pure Appl. Algebra 212 (2008), 1770--1786.
\href{https://doi.org/10.1016/j.jpaa.2007.11.010}{doi:10.1016/j.jpaa.2007.11.010}.
\end{thebibliography}
\end{document}
